\documentclass[11pt,a4paper]{article}

\usepackage[T1]{fontenc}
\usepackage[utf8]{inputenc}
\usepackage{lmodern}
\usepackage{microtype}
\usepackage[margin=2.5cm]{geometry}
\usepackage{amsmath,amssymb}
\usepackage{booktabs}
\usepackage{tabularx}
\usepackage{array}
\usepackage{enumitem}
\usepackage[hidelinks]{hyperref}

\setlist{itemsep=2pt,topsep=4pt}
\newcommand{\code}[1]{\texttt{#1}}
\newcommand{\cvxcla}{\code{cvxcla}}
\newcommand{\repo}{https://github.com/cvxgrp/cvxcla/blob/main}
\newcommand{\repofile}[1]{\href{\repo/#1}{\code{\detokenize{#1}}}}
\newcolumntype{L}{>{\raggedright\arraybackslash}X}

\title{\cvxcla{} vs.\ Bailey--López de Prado (2013):\\
  a measured comparison of two Critical Line Algorithm implementations}
\author{Thomas Schmelzer}
\date{\today}

\begin{document}
\maketitle

\begin{abstract}
  We compare the \cvxcla{} implementation of Markowitz's Critical Line
  Algorithm with the reference open-source implementation of Bailey and
  López de Prado (BLdP)~\cite{bailey2013}, as shipped in PyPortfolioOpt.
  Beyond a source-level comparison, we report three experiments. First, the
  runtime gap grows with problem size without a plateau, reaching $418\times$
  at $n = 640$; at most $1.6\times$ of it is due to the dense inverse inside
  the candidate loop. Second, the shipped BLdP code fails an exact path check on
  36 of 57 generic problems; one bug (stale blocked weights) loses real
  portfolios, while a second defect only drops $\lambda$ labels. It also fails
  on degenerate inputs that \cvxcla{} handles. Third, restricted to
  BLdP's own problem class, the gap is algorithmic: BLdP's per-step complexity
  is a factor of about $n$ above what the algorithm needs, a factor that any
  vectorised CLA recovers. \cvxcla{} can claim robustness at a constant cost,
  faster linear algebra at scale, and features outside BLdP's class, but not
  the $100\times$ headline.
\end{abstract}

\section{Introduction}

This note compares the \cvxcla{} paper (Schmelzer, \emph{The Critical Line
Algorithm as implemented in cvxcla}~\cite{cvxcla}) with Bailey and López de
Prado, ``An Open-Source Implementation of the Critical-Line Algorithm for
Portfolio Optimization''~\cite{bailey2013}.

This is not an arbitrary comparison: Bailey--López de Prado (BLdP) is the
reference that \cvxcla{} benchmarks against, and the PyPortfolioOpt~\cite{pypfopt}
\code{CLA} used in the experiments is an implementation of it. The comparison
therefore doubles as a fact-check of \cvxcla{}'s \S8 (``Relation to the
Bailey--López de Prado algorithm'').

Sections~\ref{sec:lineage} and~\ref{sec:differences} compare the two
implementations from their papers and source. Sections~\ref{sec:gap}--\ref{sec:fair}
report the three experiments of issue
\href{https://github.com/cvxgrp/cvxcla/issues/869}{\#869}, each reproducible
with a script under \code{experiments/}. Section~\ref{sec:strengths} weighs
the two papers, and Section~\ref{sec:verdict} concludes.

\section{Shared lineage}\label{sec:lineage}

Both papers compute the \emph{entire} efficient frontier exactly as a finite
list of Markowitz turning points~\cite{markowitz1956}, in Python. They share
several concrete design choices:

\begin{itemize}
  \item \textbf{First turning point.} The same greedy construction: start
    every weight at its lower bound, fill assets in descending
    expected-return order until the budget binds; the last (partially filled)
    asset is the sole free asset (BLdP \code{initAlgo}, Snippet~2;
    \cvxcla{} \code{init\_algo}, \S3).
  \item \textbf{Minimum-variance endpoint.} Both append an explicit turning
    point at $\lambda = 0$, and both explicitly note this \emph{departs from
    Niedermayer et al.}~\cite{niedermayer2010}, who keep searching for turning
    points at negative $\lambda$.
  \item \textbf{Maximum-Sharpe portfolio.} Both locate it by a
    one-dimensional search over the convex combination of the two turning
    points that bracket the discrete Sharpe maximum (BLdP: golden-section
    search, Snippets~13--14; \cvxcla{}: one-dimensional search, \S6).
  \item \textbf{Frontier interpolation.} Both reconstruct intermediate
    frontier portfolios as convex combinations of adjacent turning points.
\end{itemize}

\section{Head-to-head differences}\label{sec:differences}

Table~\ref{tab:differences} summarises the differences, verified against the
BLdP source.

\begin{table}[ht]
  \centering
  \small
  \caption{Differences between the two implementations.}\label{tab:differences}
  \begin{tabularx}{\textwidth}{>{\raggedright\arraybackslash}p{2.8cm}LL}
    \toprule
    Axis & Bailey--López de Prado (2013) & \cvxcla{} \\
    \midrule
    Covariance representation
      & Explicit dense matrix, sliced via \code{getMatrices} /
        \code{reduceMatrix} (Snippet~5)
      & Operator protocol (\code{matvec}, \code{solve\_free}, \code{cross})
        plus structured backends \\
    Free-block linear algebra
      & Explicit \code{np.linalg.inv(covarF)} at each step (Snippets~3, 7--9)
      & Single block-eliminated KKT solve via an $m \times m$ Schur
        complement; no explicit inverse \\
    Asset-entry search
      & Full inverse recomputed \emph{inside} the candidate loop
        \code{for i in b} (Snippet~3), hence $O(n)$ inverses per turning point
      & Vectorised event search; one solve per step \\
    Structured risk models
      & Not supported
      & Diagonal-plus-low-rank factor backend via the Woodbury identity (no
        $n \times n$ matrix ever formed) \\
    Constraints
      & Box bounds and a single (budget) equality only (inputs: \code{mean},
        \code{covar}, \code{lB}, \code{uB})
      & Arbitrary $m$ linear equality constraints $Aw = b$ \\
    Degeneracy handling
      & No general anti-cycling rule; a single hack (\code{mean[-1]+=1e-5})
        for all-equal means
      & Bland-style lowest-index tie-break and $\sqrt{\varepsilon}$ slope
        floor (\S4.3--4.4) \\
    \bottomrule
  \end{tabularx}
\end{table}

\subsection{The central claim checks out}

\cvxcla{}'s most load-bearing \S8/\S9.3 claim, that PyPortfolioOpt
``recomputes that full inverse once for every candidate'' in the asset-entry
search, is literally BLdP Snippet~3: \code{covarF\_inv = np.linalg.inv(covarF)}
sits \emph{inside} the \code{for i in b} candidate loop. The per-candidate loop
is what drives the runtime gap, but the measurements below show that the dense
\emph{inverse} inside it accounts for only a small part of that gap.

One nuance is worth noting: BLdP's \emph{removal} branch (Snippet~7) computes
the inverse once \emph{outside} its loop, so the per-candidate rebuild is
specific to the \emph{addition} branch, which is exactly how \cvxcla{} phrases
it.

\section{Where the gap comes from}\label{sec:gap}

The script \repofile{experiments/bldp_gap.py} re-measures the gap and breaks it
down. It uses PyPortfolioOpt~1.6.0, numpy~2.4.6 (Accelerate) and \cvxcla{}~2.6.0
on an Apple M4 Pro, with the factor-model problems of the paper's scaling
figure ($K = 10$, long-only, budget), 3 seeds $\times$ 3 repeats per size, and
each measurement in a fresh process. Ratios are medians over seeds; the spread
over seeds is within $\pm 10\%$.

\begin{itemize}
  \item \textbf{Code reading confirmed.} In PyPortfolioOpt's \code{\_solve}
    (verified by parsing the shipped source, not the paper's snippet
    numbering), the only \code{np.linalg.inv} inside a loop is the addition
    branch's \code{for i in b}. At $n = 640$ a trace performs 321 inverses per
    turning point, about $n/2$.
  \item \textbf{The gap grows with $n$, with no plateau.} BLdP/\cvxcla{} is
    $2.8\times$ at $n = 20$, $16\times$ at 80, $146\times$ at 320 and
    $\mathbf{418\times}$ at 640, against the $\approx 340\times$ reported. Over
    the last doubling, BLdP scales like $n^{3.7}$ and \cvxcla{} like $n^{2.1}$.
  \item \textbf{The gap breaks into three factors}, each from changing one
    thing at a time (Table~\ref{tab:gap}).
\end{itemize}

\begin{table}[ht]
  \centering
  \small
  \caption{Factors of the runtime gap (track~1). Each column changes one thing
    at a time; their product is the total.}\label{tab:gap}
  \begin{tabular}{rrrrr}
    \toprule
    $n$ & \shortstack[r]{inverse $\to$ bordered\\update, loop kept}
        & \shortstack[r]{Python loop $\to$\\vectorised search}
        & \shortstack[r]{minimal vectorised\\CLA $\to$ \cvxcla{}}
        & total \\
    \midrule
    20  & $0.9\times$ & $12\times$  & $0.3\times$ & $2.8\times$ \\
    80  & $1.0\times$ & $37\times$  & $0.4\times$ & $16\times$  \\
    320 & $1.3\times$ & $86\times$  & $1.3\times$ & $146\times$ \\
    640 & $1.6\times$ & $141\times$ & $1.9\times$ & $418\times$ \\
    \bottomrule
  \end{tabular}
\end{table}

Replacing each candidate's fresh $O(n_F^3)$ inverse with an $O(n_F^2)$
bordered update, while keeping the loop, buys at most $1.6\times$. Most of the
gap comes from evaluating $O(n)$ candidates one at a time. (Section~\ref{sec:fair}
shows that this cost is the per-candidate \emph{formula}, not the Python loop:
the loop itself costs $1.2$--$1.5\times$.) Each evaluation re-slices
\code{covarF}/\code{covarFB} and re-runs the $\lambda$ formula, so it costs
$O(n \cdot n_F)$ with or without an inverse. A phase profile of the shipped
code agrees: the addition-branch inverses take 8\% of the trace at $n = 20$
and 41\% at $n = 640$. The slicing takes 27--40\%, and the $\lambda$ formula
plus the removal branch most of the rest. The last factor is \cvxcla{}'s own
cost relative to a $\sim$100-line vectorised CLA: \cvxcla{} is $3\times$
\emph{slower} at $n = 20$ (validation and guards), and $1.9\times$ faster at
$n = 640$ (a block-eliminated solve instead of an explicit inverse).

\begin{itemize}
  \item \textbf{Maintaining the inverse pays off at scale, in both codes.}
    Updating the free-block inverse by rank-one steps instead of rebuilding it
    makes the minimal CLA $3.9\times$ faster at $n = 640$, and makes \cvxcla{}
    $2.5\times$ faster with \code{IncrementalDenseCovariance} instead of the
    default dense backend (no gain below $n = 80$). With matching linear
    algebra, \cvxcla{}'s incremental backend is $3\times$ slower than the
    minimal CLA at $n = 20$ but $1.25\times$ \emph{faster} at $n = 640$
    (204\,ms against 247\,ms), and it scales better ($n^{1.7}$ against
    $n^{2.3}$ over the last doubling). These numbers come from a separate run of
    the four fast variants; the BLdP columns of Table~\ref{tab:gap} were not
    re-timed with it.
  \item \textbf{BLdP can cycle.} Its re-entry test \code{lam < self.ls[-1]}
    has no tolerance and its loop has no iteration cap. With bordered updates,
    which change $\lambda$ by about $10^{-15}$, an asset alternately leaves and
    re-enters at the same $\lambda$ forever: 5 of the 18 runs of that variant
    never terminate. The shipped code finished on all 18 problems, but only by
    the luck of its rounding. (While checking the vectorised variants against
    \cvxcla{}, the experiment also found that \code{experiments/inverse\_cla.py}
    used an absolute $10^{-5}$ tie window on $\lambda$ that stepped past
    turning points from $n = 160$ on, with weights down to $-5 \times 10^{-6}$
    under a zero bound. It now uses $10^{-9}$.)
  \item \textbf{Not measured:} other CPUs or BLAS builds (\code{--threads} is
    available but was not swept), $n > 640$, other problem families, memory,
    or the original BLdP code as published rather than PyPortfolioOpt's
    descendant of it.
\end{itemize}

\section{Do they compute the same frontier?}\label{sec:agree}

The script \repofile{experiments/bldp_agree.py} feeds the same \code{mean},
\code{covariance}, \code{lB} and \code{uB} to PyPortfolioOpt~1.6.0's \code{CLA}
and to \cvxcla{}~2.7.0. It checks each trace on its own against
CVXPY~1.9.3~\cite{cvxpy} / Clarabel~0.11.1~\cite{clarabel} at tolerance
$10^{-12}$, at every turning point and every segment midpoint, and compares the
two paths $w(\lambda)$ exactly. Both use the same $\lambda$ convention. There
are 90 problems in eight families over 3 seeds: generic, shorting, all-equal
means, tied means, duplicated assets, rank-deficient sample covariances, a
conditioning sweep from $10^{2}$ to $10^{16}$, and bounds that bind exactly.
``Exact'' means feasible to $10^{-9}$, $\lambda$ running from $\infty$ to 0,
and a relative objective gap below $10^{-8}$ everywhere.

\subsection{On generic problems the shipped code is often wrong}

\cvxcla{} is exact on all 57 feasible non-degenerate problems except two at
$\kappa \ge 10^{14}$ (Section~\ref{sec:conditioning}). The shipped BLdP code is
exact on only 21 of them. The rest fail the check, with objective gaps up to
$3 \times 10^{-4}$ and paths up to $0.35$ apart in max-norm, and no error is
raised. Even plain long-only problems fail on 4 of 12 ($n = 10$ to 100), all
from the first defect below. The check compares the two paths as functions of
$\lambda$, so it counts two different defects, and only the first one loses
portfolios:

\begin{enumerate}
  \item \textbf{Stale blocked weights.} When a free asset hits a bound,
    \code{\_solve} sets the working vector \code{w[i] = bound}.
    \code{\_get\_matrices} then reads the blocked weights from
    \code{self.w[-1]}, the \emph{previous} turning point. The new point is
    solved with asset $i$ at its old free value, so its budget is off by
    exactly $\text{bound} - w_{\text{prev},i}$ (e.g.\ $-4.1 \times 10^{-2}$).
    \code{\_purge\_num\_err} deletes it as infeasible, and a true vertex goes
    with it. Sometimes the trace recovers instead by removing, re-adding and
    removing the asset at the same $\lambda$. That leaves three copies of one
    point, the same leave-and-re-enter pattern that Section~\ref{sec:gap} saw
    cycle forever under bordered updates. They do not cause that cycle:
    Section~\ref{sec:fair} found the bordered variant cycling on the same 9 of
    30 problems with or without both fixes. Reading the blocked weights from
    the working vector (\code{bldp-fix-w}, a two-method subclass) makes 52 of
    the 57 exact. This bug is user-visible: PyPortfolioOpt issue
    \href{https://github.com/PyPortfolio/PyPortfolioOpt/issues/750}{\#750}
    reports it independently, with \code{max\_sharpe} silently suboptimal in
    22 of 300 random problems.
  \item \textbf{A zero-tolerance purge, which loses only $\lambda$ labels.}
    \code{\_purge\_excess} deletes a turning point whenever a later one has a
    strictly higher $\mu^\top w$. Where $\mu^\top w$ is flat, rounding makes
    it rise by about $10^{-14}$. Along an exact trace, though, $\mu^\top w$ is
    flat only where the weights are constant, since
    $\mathrm{d}(\mu^\top w)/\mathrm{d}\lambda = \mu^\top (\mathrm{d}w/\mathrm{d}\lambda)$
    vanishes only when $\mathrm{d}w/\mathrm{d}\lambda = 0$. So every point the
    purge deletes is a copy of a kept turning point. With the first bug fixed,
    it deletes 34 points in 17 of the 57 problems, and every one matches a kept
    point to within $1.1 \times 10^{-11}$. No portfolio is lost, and
    \code{max\_sharpe}, which interpolates by weight and never reads
    $\lambda$, is unchanged. What is lost is the $\lambda$ at which the
    weights stop being constant; the check here reads the path in $\lambda$,
    so it scores this as an error. This defect was first reported as deleting
    true vertices; the reporter of \#750 showed otherwise, and the correction
    is posted there. With a $10^{-12}$ relative tolerance as well
    (\code{bldp-fix-wp}), 56 of 57 are exact as $\lambda$-paths with gaps
    $\le 7 \times 10^{-13}$. The remaining case cycles at
    $\kappa = 10^{16}$.
\end{enumerate}

With both fixes, BLdP and \cvxcla{} trace the same path to $10^{-8}$ on every
one of these problems except the two below. Only the first fix is needed for
the same set of portfolios. The defects are in the implementation, not in the
algorithm.

\subsection{Degenerate inputs break BLdP; \cvxcla{} handles every one}

\cvxcla{} is exact on all 30 problems in the all-equal, tied, duplicated and
rank-deficient families. BLdP, with or without the two fixes:

\begin{itemize}
  \item \emph{All-equal means:} raises \code{TypeError} on all 6.
    PyPortfolioOpt ships the paper's \code{mean[-1] += 1e-5} hack commented
    out. With it switched back on, BLdP still raises for $n = 20$, and for
    $n = 5$ traces the frontier of a perturbed problem whose $\lambda$ is
    scaled by $1/10^{-5}$ (gap $\approx 4 \times 10^{-5}$). The true frontier
    is the single minimum-variance point.
  \item \emph{Ties in expected returns:} raises \code{TypeError} on 7 of 9.
    When the event-search denominator $c$ is exactly 0,
    \code{\_compute\_lambda} returns \code{None}, and the comparison
    \code{lam < self.ls[-1]} fails on it. Once, rounding kept $c$ off zero and
    the trace ran silently down the wrong path: \code{\_init\_algo} starts at
    one vertex of the tied face and never frees the other tied assets, giving a
    gap of $0.25$. \cvxcla{} starts at the face minimiser.
  \item \emph{Duplicated assets / rank-deficient $\Sigma$:} \code{LinAlgError}
    on 8 of 9 duplicated problems and \code{TypeError} on the 3 problems with
    $T < n/2$. Worse, when an exact duplicate of the top asset enters later,
    \code{np.linalg.inv} of the singular free block (condition number
    $1.5 \times 10^{16}$) does not raise. Every later point violates the
    budget, \code{\_purge\_num\_err} deletes them all, and the frontier
    silently stops at $\lambda = 11.4$ instead of 0. With $T = 30 < n = 40$ the
    fixed code is exact on 2 of 3 problems and goes wrong near $\lambda = 0$ on
    the third (gap $2 \times 10^{-6}$).
  \item \emph{Infeasible bounds} ($\sum_i u_i < 1$): \code{IndexError} from
    \code{\_init\_algo} walking off the asset list. \cvxcla{} raises
    \code{InfeasibleProblemError}.
\end{itemize}

\subsection{Which degrades first with conditioning?}\label{sec:conditioning}

With both fixes, BLdP's explicit inverse holds up surprisingly well: gaps
$\le 4 \times 10^{-13}$ up to $\kappa = 10^{14}$, and one cycle at
$10^{16}$. \cvxcla{} stays exact up to $10^{12}$. At $\kappa = 10^{14}$ and
$10^{16}$ it has one problem each with a gap of $5 \times 10^{-8}$ and
$8 \times 10^{-8}$. That is on a segment at $\lambda \approx 10^{-10}$, where
its weights are $2 \times 10^{-3}$ off the QP's along a near-null direction of
$\Sigma$, while BLdP-fixed and the QP agree to $10^{-14}$. As shipped, BLdP is
wrong from $\kappa = 10^{2}$ because of the two defects.

\paragraph{Not measured.} The 2013 paper's own code as published rather than
PyPortfolioOpt's descendant (whether defects 1 and 2 are in the original is not
checked here), the Markowitz--Todd 10-asset example~\cite{markowitz2000},
$n > 100$, and constraint classes beyond a budget and box bounds (BLdP
supports nothing else).

\section{Is the comparison apples to apples?}\label{sec:fair}

The script \repofile{experiments/bldp_fair.py} re-times the gap with every rung
restricted to BLdP's own problem class: dense $\Sigma$, box bounds and one
budget row. It uses the $K = 10$ factor problems of Section~\ref{sec:gap},
$n = 20$ to 320, 3 seeds $\times$ 3 repeats, each in a fresh process, on the
same M4 Pro, Python~3.12 and numpy~2.4.6 stack. It asks which parts of the gap
are \cvxcla{}'s to claim.

\subsection{Correcting BLdP does not shrink the gap}

With both fixes of Section~\ref{sec:agree}, BLdP is 4--9\% \emph{slower}
than as shipped, because it keeps the turning points the purges used to delete.
The gap against a correct BLdP is $3.5\times$ at $n = 20$ and $150\times$ at
$n = 320$, against $3.1\times$ and $144\times$ for the shipped code. The ratios
of Section~\ref{sec:gap} therefore stand, but they were measured against code
that returns a wrong frontier on many of these problems (exact distance to
\cvxcla{}'s path up to $10^{-2}$).

\subsection{The gap is algorithmic, not Python}

The $12$--$141\times$ ``vectorised event search'' rung of Table~\ref{tab:gap}
changed two things at once. A new rung, \code{loop} (an option of
\code{InverseCLA}), keeps the reduced-cost formula of the vectorised code but
visits the assets one at a time in a Python loop. That splits the rung in two
(Table~\ref{tab:fair}).

\begin{table}[ht]
  \centering
  \small
  \caption{Factors of the gap on BLdP's own problem class (track~3).}\label{tab:fair}
  \begin{tabular}{rrrrrr}
    \toprule
    $n$ & bug fixes & formula & interpreter & \cvxcla{} vs minimal & total \\
    \midrule
    20  & $0.91\times$ & $11\times$ & $1.22\times$ & $0.23\times$ & $3.1\times$ \\
    80  & $0.92\times$ & $27\times$ & $1.50\times$ & $0.44\times$ & $16\times$  \\
    320 & $0.96\times$ & $83\times$ & $1.45\times$ & $1.24\times$ & $144\times$ \\
    \bottomrule
  \end{tabular}
\end{table}

The \emph{formula} factor grows like $n^{0.96}$. BLdP re-derives each
candidate's $\lambda$ from its own bordered system, re-slicing $\Sigma$ and
recomputing the shared products each time: $O(n)$ candidates $\times$
$O(n_F^3 + n \cdot n_F)$ work per turning point. Reading every candidate's
$\lambda$ off one solve of the current segment ($-\gamma_i/\delta_i$) costs
$O(n)$ per candidate. That is a change in complexity, and it accounts for
almost all of the gap. It includes the per-candidate inverse, which
Section~\ref{sec:gap} put at $\le 1.6\times$. The \emph{interpreter} factor,
the same flops dispatched once per asset instead of once per step, is flat at
$1.2$--$1.5\times$ ($n^{-0.1}$). A Python loop is not what makes BLdP slow.

\subsection{What \cvxcla{}'s architecture is entitled to claim}

None of the formula factor is specific to \cvxcla{}: a $\sim$100-line
vectorised CLA (\code{vec-fresh}) has it too. \cvxcla{}'s own contribution is
the last column of Table~\ref{tab:fair}. It is $4.3\times$ \emph{slower} than
the minimal code at $n = 20$ and $1.24\times$ faster at $n = 320$, crossing
over near $n = 250$. A cProfile of \cvxcla{} shows where its extra time goes:

\begin{itemize}
  \item \emph{Guards} take 48\% of its trace at $n = 20$ (1.3\,ms) and 25\% at
    $n = 320$ (26\,ms). These are the feasibility checks and projection,
    held-bound checks, active-row guards, the multipliers' round-off level,
    the mean shift and tied-start handling.
  \item \emph{Generality} (the inequality-row event search, idle here) takes
    $\le 1\%$.
  \item The \emph{core} (75\% at $n = 320$, 79\,ms) is faster than the minimal
    code's whole trace (131\,ms): a block-eliminated Cholesky solve rather
    than an explicit inverse.
\end{itemize}

So on BLdP's problem class, supporting arbitrary $Aw = b$ and $Gw \le h$ costs
\cvxcla{} almost nothing. Its robustness costs a constant share, and its linear
algebra wins at scale.

\subsection{The environment does not explain the gap}

Both codes ran on the same machine and stack, so hardware cancels in every
ratio. Single-threaded BLAS changes no timing by more than 7\%, presumably
because the blocks are too small for threads to help. The shipped BLdP code at
$n = 160$ takes about 1.43\,s on Python 3.12, 3.13 and 3.14 alike
(numpy~2.4.6). It takes \textbf{3.83\,s} under numpy~1.26.4, $2.6\times$
slower. A newer numpy \emph{narrows} the gap, so the ratios above are, if
anything, kind to BLdP compared with the stack it was written for.

\subsection{What a scaling comparison may conclude}

BLdP set out to make the CLA available and verifiably correct at $n = 10$. It
never claimed speed, and at $n = 20$, the smallest size timed here, the gap is
about $3\times$. The measured gap supports a narrow statement: \emph{the
reference implementation, as shipped and widely used, has a per-step
complexity a factor of about $n$ above what the algorithm needs.} That is a
property of the code's formulation, not of the CLA, not of Python, and mostly
not of \cvxcla{}'s architecture: the formula factor ($11$--$83\times$) belongs
to any vectorised CLA. \cvxcla{} can claim the last column, robustness at a
constant cost, and the features outside BLdP's class (general constraints,
factor backends). It cannot claim the $100\times$ headline over BLdP.

\paragraph{Not measured.} Stacks older than numpy~1.26 or Python~3.12
(PyPortfolioOpt~1.6.0 and \cvxcla{} do not install there, so the 2013
environment is out of reach), other CPUs or BLAS builds, $n > 320$ for this
split, and memory. The profile fractions are taken under cProfile, which
inflates Python-call-heavy code, so the guards' share is an upper estimate.

\section{Where each paper is stronger}\label{sec:strengths}

\subsection{Bailey--López de Prado, on its own terms}

\begin{itemize}
  \item \textbf{Availability.} In 2013 it filled a genuine gap: it was the
    first documented open-source CLA in a scientific language (prior art was
    the Markowitz--Todd VBA-Excel implementation~\cite{markowitz2000}). Its
    code became PyPortfolioOpt's \code{CLA} and is widely cited and used. That
    impact is real, whatever the code's quality.
  \item \textbf{A reference answer.} It reproduces the canonical
    Markowitz--Todd 10-asset example to 15 significant figures. One example is
    a narrow check, though: Section~\ref{sec:agree} found the code wrong on 4
    of 12 ordinary long-only problems from $n = 10$ up.
\end{itemize}

It has no pedagogical edge. As code to learn from, it shows how a parametric
QP should \emph{not} be solved:

\begin{itemize}
  \item It re-derives every candidate event from its own bordered system,
    about $n$ times the work per turning point that the algorithm needs
    (Section~\ref{sec:fair}).
  \item It compares $\lambda$ values with no tolerance, so it can cycle
    (Section~\ref{sec:gap}).
  \item It hides infeasible turning points by deleting them, which turns the
    stale-weight bug into a silently wrong frontier (Section~\ref{sec:agree}).
\end{itemize}

\subsection{\cvxcla{}}

\begin{itemize}
  \item \textbf{A better entry into the field.} \cvxcla{} states each segment
    as the KKT system of the parametric QP: the optimality conditions on the
    free set, multipliers for the active constraints, and events where a
    weight hits a bound or a multiplier changes sign. That is the language of
    the parametric-QP and homotopy literature. The payoff is transferable
    understanding. The same \code{ParametricProblem} / \code{trace} engine runs
    the LARS/LASSO homotopy, with the same four events under different
    names~\cite{schmelzer2026}. BLdP's bespoke $\lambda$ formulas, re-derived
    per candidate, teach none of that structure.
  \item \textbf{Scale demonstrated.} BLdP only ever runs its $n = 10$ toy
    example. \cvxcla{} shows $n = 640$ benchmarks, an $n = 20{,}000$ factor
    stress test, and a real $n = 494$ S\&P~500 frontier.
  \item \textbf{Robustness.} BLdP never discusses conditioning, cycling, or
    tolerances and has no general anti-cycling rule. \cvxcla{} hardens the
    event logic \emph{and} honestly documents where it still fails (\S9.6).
  \item \textbf{Representation-agnosticism.} The covariance-operator
    abstraction (and the resulting sub-quadratic factor backend) is the one
    element not present in BLdP, or, per \cvxcla{}, in any standard CLA
    implementation.
  \item \textbf{Independent exactness check.} \cvxcla{} validates against a
    general convex QP solver (CVXPY/Clarabel) at solver precision, rather than
    against a single known reference answer.
\end{itemize}

\section{Ranking}\label{sec:ranking}

\cvxcla{} ranks clearly ahead, and the deciding axis is correctness, not
speed. All evidence comes from this repository's own experiments, written by
\cvxcla{}'s author, and none of it has been reproduced independently.

\subsection{Axis by axis}

\begin{table}[ht]
  \centering
  \small
  \caption{Ranking by axis.}\label{tab:ranking}
  \begin{tabularx}{\textwidth}{>{\raggedright\arraybackslash}p{2.6cm}
      >{\raggedright\arraybackslash}p{1.9cm}L}
    \toprule
    Axis & Ahead & Evidence \\
    \midrule
    Correctness & \cvxcla{}
      & Shipped BLdP loses real turning points through stale blocked weights
        (silently suboptimal \code{max\_sharpe} in \#750) and fails on tied,
        duplicated and rank-deficient inputs; \cvxcla{} is exact on all
        but two problems at $\kappa \ge 10^{14}$ (Section~\ref{sec:agree}). \\
    Robustness & \cvxcla{}
      & BLdP has no tolerances, can cycle, and deletes infeasible points;
        \cvxcla{} uses tie-breaks, tolerances and raises on infeasibility. \\
    Scope & \cvxcla{}
      & BLdP: box bounds and a budget only. \cvxcla{}: general $Aw = b$,
        $Gw \le h$, leverage caps, factor backends, the LASSO path. \\
    Speed at scale & \cvxcla{}
      & $3\times$ at $n = 20$ to $418\times$ at $n = 640$, mostly a factor of
        about $n$ that any vectorised CLA recovers; \cvxcla{}'s own share is
        $1.2$--$1.9\times$ (Sections~\ref{sec:gap},~\ref{sec:fair}). \\
    Small-$n$ overhead & neither
      & \cvxcla{} is $4\times$ slower than a minimal vectorised CLA at
        $n = 20$ because of its guards, but still faster than BLdP. \\
    Extreme conditioning & BLdP (fixed)
      & With both fixes, BLdP's explicit inverse matches the QP slightly
        better at $\kappa \ge 10^{14}$ (Section~\ref{sec:conditioning}). \\
    History and reach & BLdP
      & First open-source CLA in Python; widely used through PyPortfolioOpt. \\
    Simplicity & BLdP
      & Shorter, numpy only; but a poor model to learn from, since its
        per-candidate formulas hide the KKT structure. \\
    \bottomrule
  \end{tabularx}
\end{table}

\subsection{Overall}

\begin{enumerate}
  \item \textbf{\cvxcla{}}, for any real use.
  \item \textbf{BLdP with the stale-weight fix}: correct on generic problems, but
    still about $n$ times too slow and fragile on degenerate inputs.
  \item \textbf{BLdP as shipped}: not recommended; its results are wrong too
    often, with no error raised.
\end{enumerate}

The ranking does not credit \cvxcla{}'s architecture with the $100\times$-plus
speed headline. Most of that comes from writing the CLA properly, not from
\cvxcla{} specifically: a minimal correct vectorised CLA would also beat BLdP
by a wide margin. \cvxcla{} adds robustness, scope and somewhat faster linear
algebra on top of that.

Two checks would make the ranking more convincing: running the original 2013
code rather than PyPortfolioOpt's descendant of it, and having someone outside
the project repeat the experiments, as issue
\href{https://github.com/cvxgrp/cvxcla/issues/869}{\#869} invites.

\section{Verdict}\label{sec:verdict}

BLdP's lasting contribution is \textbf{availability}. It established the
open-source baseline, and its code is in wide use through PyPortfolioOpt. It
is not a model implementation:

\begin{itemize}
  \item its event search does about $n$ times the necessary work per step;
  \item it has no tolerances;
  \item it repairs infeasible points by deleting them.
\end{itemize}

The good practice is the reverse: write down the KKT system of the parametric
QP, solve the current segment once and read every event off it, compare with
tolerances, and validate against an independent solver. That is also the more
helpful entry into the field, because it is how the parametric-QP and homotopy
literature, LARS/LASSO included, presents path following. \cvxcla{} is an
\textbf{engineering / performance paper} that rebuilds the method on those
principles and quantifies the payoff: operator abstraction, block elimination,
structured backends and hardened degeneracy handling.

\cvxcla{}'s characterisation of BLdP is accurate. If anything, it is too narrow
about the cause of the gap. It blames the inverse inside the candidate loop,
which accounts for at most $1.6\times$, where Sections~\ref{sec:gap}
and~\ref{sec:fair} locate the gap in the per-candidate formulation as a whole.
\S8 could now also note that the reference code returns wrong frontiers on
ordinary problems (Section~\ref{sec:agree}).

\begin{thebibliography}{9}

\bibitem{bailey2013}
D.~H. Bailey and M.~López de Prado.
\newblock An open-source implementation of the critical-line algorithm for
  portfolio optimization.
\newblock \emph{Algorithms}, 6(1):169--196, 2013.
\newblock SSRN abstract 2197616.

\bibitem{cvxcla}
T.~Schmelzer.
\newblock The Critical Line Algorithm as implemented in cvxcla.
\newblock Software and manuscript, \url{https://github.com/cvxgrp/cvxcla}.

\bibitem{schmelzer2026}
T.~Schmelzer and T.~Hastie.
\newblock The Critical Line Algorithm and the constrained LASSO: one curve,
  two literatures.
\newblock arXiv:2609.25704 [stat.ME], 2026.

\bibitem{markowitz1956}
H.~Markowitz.
\newblock The optimization of a quadratic function subject to linear
  constraints.
\newblock \emph{Naval Research Logistics Quarterly}, 3(1--2):111--133, 1956.

\bibitem{markowitz2000}
H.~M. Markowitz and G.~P. Todd.
\newblock \emph{Mean-Variance Analysis in Portfolio Choice and Capital
  Markets}.
\newblock Wiley, 2000.

\bibitem{niedermayer2010}
A.~Niedermayer and D.~Niedermayer.
\newblock Applying Markowitz's critical line algorithm.
\newblock In J.~B. Guerard, editor, \emph{Handbook of Portfolio Construction},
  pages 383--400. Springer, 2010.

\bibitem{pypfopt}
R.~A. Martin.
\newblock PyPortfolioOpt: portfolio optimization in Python.
\newblock \emph{Journal of Open Source Software}, 6(61):3066, 2021.

\bibitem{cvxpy}
S.~Diamond and S.~Boyd.
\newblock CVXPY: A Python-embedded modeling language for convex optimization.
\newblock \emph{Journal of Machine Learning Research}, 17(83):1--5, 2016.

\bibitem{clarabel}
P.~J. Goulart and Y.~Chen.
\newblock Clarabel: An interior-point solver for conic programs with
  quadratic objectives.
\newblock arXiv:2405.12762, 2024.

\end{thebibliography}

\end{document}
